\section{Special Functions}
\renewcommand{\theequation}{\thesection.\arabic{equation}}
\setcounter{equation}{0}
This appendix collects special functions used in the note.


\subsection{Chebyshev Polynomials}\label{subsec:Cheby}
The \textit{first kind of Chebyshev polynomial} over the real interval $[-1;1]$ is defined as 
\begin{equation}
    \textbf{T}_j(x)\coloneqq\cos (j \arccos(x)).
    \label{eqn:FirstChebyDef}
\end{equation}
Let $\theta=\arccos(x)$, note that
\begin{equation}
    \cos(j\theta)=\sum_{\substack{0 \le k \le j \\ k \text{ is even}}}\binom{j}{k}i^k\cos^{j-k}(\theta)(1-\cos^2(\theta))^{k/2},
    \label{eqn:ExpanCosjtheta}
\end{equation}
 we can see Eqn.~(\ref{eqn:FirstChebyDef}) can be written into
 \begin{equation}
    \textbf{T}_j(x)=\sum_{\substack{0 \le k \le j \\ k \text{ is even}}}\binom{j}{k}i^k x^{j-k}(1-x^2)^{k/2},
    \label{eqn:FirtChebExpansionForm}
 \end{equation}
 with which we can see $\textbf{T}_j(x)$ is indeed a polynomial. By the way, to prove Eqn.~(\ref{eqn:ExpanCosjtheta}), we need to realize, according to Euler's equation, we have 
 \begin{equation}
    \cos(j\theta)=\Re[(\cos\theta+i\sin\theta)^j].
 \end{equation}

 Now, let's consider the power series $\sum_{j=0}^{\infty}\textbf{T}_j(x)y^j$ generated by the Chebyshev polynomial. Assumming $|y|<1$, we have 
 \begin{equation}
    \begin{split}
        \sum_{j=0}^{\infty}\textbf{T}_j(x)y^j &=\sum_{j=0}^{\infty}\cos(j\theta)y^j=\sum_{j=0}^{\infty}\frac{e^{ij\theta+}+e^{-ij\theta}}{2}y^j\\
        &=\frac{\frac{1}{1-ye^{i\theta}}+\frac{1}{1-ye^{-i\theta}}}{2}=\frac{1-yx}{1+y^2-2yx}
    \end{split}
 \end{equation}
 Sometimes, we use rescaled first Chebyshev polynomial by a factor of 1/2:
 \begin{equation}
    \tilde{\textbf{T}}_j(x)=
    \begin{cases}
        \textbf{T}_j(x), \quad j\ge 1, \\
        \frac{1}{2}\textbf{T}_0(x), \quad j=0.
    \end{cases}
 \end{equation}
In this case, the alternative generative function is
\begin{equation}
    \sum_{j=0}^{\infty}\tilde{\textbf{T}}_j(x)y^j=\frac{1}{2}+\sum_{j=1}^{\infty}\textbf{T}_j(x)y^j=\frac{1-y^2}{2(1+y^2-2yx)}.
\end{equation}
The \textit{Chebyshev polynomial of the second kind} is defined as
\begin{equation}
    \textbf{U}_j(x)\coloneqq\frac{\sin((j+1)\arccos(x))}{\sin(\arccos(x))},
\end{equation}
which has the generating function for $|y|<1$
\begin{equation}
    \sum_{j=0}^{\infty}\textbf{U}_j(x)y^j=\frac{1}{1+y^2-2yx}.
\end{equation}
The two kinds of Chebyshev polynomial are related as the solutions of the Pell equation
\begin{equation}
    \textbf{T}^2_j(x)-(x^2-1)\textbf{U}_{j-1}^2(x)=1.
\end{equation}
The orthogonal relation between the Chebyshev polynomials is
\begin{equation}
    \int_{-1}^{1}\textbf{T}_j(x)\textbf{T}_k(x)\frac{dx}{\sqrt{1-x^2}}=
    \begin{cases}
        0, \quad j\neq k,\\
        \frac{\pi}{2}, \quad j=k\neq0, \\
        \pi, \quad j=k=0.
    \end{cases}
\end{equation}
Therefore, if a function has a Chebyshev expasion
\begin{equation}
    f(x)=\sum_{j=0}^{\infty}\beta_j\textbf{T}_j(x),
\end{equation}
the expansion coefficients must be uniquely determined by
\begin{equation}
    \beta_j=
    \begin{cases}
        \frac{2}{\pi}\int_{-1}^{1}\textbf{T}_j(x)f(x)\frac{dx}{\sqrt{1-x^2}}, \quad j\ge1, \\
        \frac{1}{\pi}\int_{-1}^{1}\textbf{T}_0(x)f(x)\frac{dx}{\sqrt{1-x^2}}, \quad j=0.
    \end{cases}
\end{equation}


\subsection{Error Function} \label{sub:ErrorFun}
The following content are mainly from \href{https://en.wikipedia.org/wiki/Error_function}{Wikipedia}. In mathematics, the \textbf{error function} (also called the \textbf{Gauss error function}), often denoted by \textbf{erf}, is a function $\text{erf}:\mathbb{C}\rightarrow\mathbb{C}$ defined as:
\begin{equation}
    \operatorname{erf}(z)=\frac{2}{\sqrt{\pi}}\int_{0}^{\pi}e^{-t^2}dt.
\end{equation}
In general, this integral is a complex contour integral which is path-independent because $\exp(-t^2)$ is holomorphic on the whole complex plane $\mathbb{C}$. A holomorphic function is a complex-valued function of one or more complex variables that is complex differentiable in a neighbourhood of each point in a domain in complex coordinates space $\mathbb{C}^n$. In many applications, the function argument is a real number, in which case the function value is also real. This nonelementary integral is a sigmoid function that occurs oftern in probability, statistics, and partial differential equations.
\begin{figure}
    \centering
    \includegraphics[width=0.4\textwidth]{Error_Function.svg.png}
    \caption{Error Function}
    \label{fig:ErrorFunction}
\end{figure}

In statistics, for non-negative real values of $x$, the error function has the following interpretation: for a real random variable $Y$ which is normally distributed with mean 0 and standard deviation $\frac{1}{\sqrt{2}}$, $\operatorname{erf}(x)$ is the probability that $Y$ falls in the range $[-x, x]$.

Two closely related functions are the \textbf{complementary error function} $\operatorname{erfc}:\mathbb{C}\rightarrow\mathbb{C}$ is defined as
\begin{equation}
    \operatorname{erfc}(z)=1-\operatorname{erf}(z),
\end{equation}
and the \textbf{imaginary error function} $\operatorname{erfi}:\mathbb{C}\rightarrow\mathbb{C}$ is defined as 
\begin{equation}
    \operatorname{erfi}(z)=-i\operatorname{erf}(iz).
\end{equation}
Its derivative is 
\begin{equation}
    \frac{d}{dz}\operatorname{erf}(z)=\frac{2}{\sqrt{\pi}}e^{-z^2}.
\end{equation}
Its antiderivative is
\begin{equation}
    \int\operatorname{erf}(z)dz=z\operatorname{erf}(z)+\frac{e^{-z^2}}{\sqrt{\pi}}+C.
\end{equation}
Its Taylor series expansion is
\begin{equation}
    \operatorname{erf}(z)=\frac{2}{\sqrt{\pi}}\sum_{n=0}^{\infty}\frac{(-1)^n}{2n+1}\frac{z^{2n+1}}{n!}.
\end{equation}